% =====================================================================
\section{Stress tests}\label{sec:stress}
% =====================================================================

The kernel fit of \S\ref{sec:results} is stress-tested in five
independent ways: (i) bin bootstrap, (ii) $q^{2}$ region splits, (iii)
alternative single- and two-Wilson-coefficient comparison models,
(iv) a charm-loop nuisance proxy, and (v) a 20-sample BSZ
form-factor Monte Carlo. The kernel shape is held fixed throughout;
only $A$ is fitted.

\paragraph{Mode caveat.} All stress tests in this section use the
linearised Mode-B response (Eq.~\ref{eq:modeB}). The headline non-linear
analysis on LHCb 2025 (\S\ref{subsec:joint_fit}) gives different
absolute $\chi^{2}$ values, smaller best-fit amplitudes, and a
$\Delta\mathrm{AIC}$ of opposite sign vs $\mathrm{FREE\_C9}$ (see
Table~\ref{tab:joint_fit_modes}). What the Mode-B stress tests
\emph{can} establish is qualitative: sign-stability of $A$ under
perturbations, region-by-region behaviour, and the relative
performance of competing $q^{2}$-shapes. They cannot rescue the
non-linear $\Delta\mathrm{AIC}$ result. Numbers below are quoted
under Mode B and should be read accordingly.

\subsection{Bin bootstrap}

A non-parametric bootstrap is run with $N=500$ trials. Each trial
resamples 8 $q^{2}$ bins with replacement from the canonical grid;
all four observables in a chosen bin enter together (bins are the
unit of resampling, not individual observable rows). Loss uses
diagonal stat$\oplus$syst because resampling with replacement makes
the published covariance singular for duplicated rows.

\begin{table}[h]
\centering
\small
\begin{tabular}{lc}
\toprule
statistic & value \\
\midrule
mean $A$         & $+1.592$ \\
median $A$       & $+1.594$ \\
std $A$          & $0.203$ \\
$90\,\%$ CI      & $[+1.356,\ +1.852]$ \\
fraction $A<0$   & $0.002$ \\
\bottomrule
\end{tabular}
\caption{Bootstrap statistics for the shared kernel amplitude on
the LHCb 2025 four-observable joint fit, $N=500$ trials. Sign-stable
to $99.8\,\%$; 90\,\% CI within $\pm 15\,\%$ of the central value.}
\label{tab:bootstrap}
\end{table}

\subsection{Region splits}

The same fit is repeated independently on three $q^{2}$ sub-windows.

\begin{table}[h]
\centering
\small
\begin{tabular}{lcrrrrrr}
\toprule
region & $q^{2}$ range (GeV$^{2}$) & $n_{\mathrm{bins}}$ & $n$ &
   FREE\_C9 $\chi^{2}$ & FREE $\Delta C_{9}$ &
   VFD $A$ & $\Delta\mathrm{AIC}$ \\
\midrule
low\_q2     & $[0.06,\,4.0]$  & 3 & 12 & 20.86 & $-1.13$ & $+1.40$ & $-0.09$ \\
central\_q2 & $[4.0,\,8.0]$   & 2 &  8 &  5.43 & $-1.56$ & $+1.82$ & $-0.33$ \\
high\_q2    & $[11.0,\,19.0]$ & 3 & 12 & 10.04 & $-1.17$ & $+1.30$ & $-0.69$ \\
\bottomrule
\end{tabular}
\caption{Region-split fits on the LHCb 2025 four-observable joint
data with the kernel held fixed. VFD beats FREE\_C9 in every region.
Per-region amplitudes are $A\in[1.30, 1.82]$ (sign-uniform), with
$\Delta C_{9}^{\mathrm{eff}}\in[-1.57, -1.12]$ (sign-uniform).}
\label{tab:regions}
\end{table}

The amplitude varies in a 40\,\% band across regions but stays
sign-uniform; the high-$q^{2}$ window — conventionally the regime
where charm-loop / long-distance contributions are largest and where
SM theory is least reliable — is the region where VFD's AIC advantage
is largest, not smallest.

\subsection{Alternative single- and two-Wilson-coefficient models}

We compare against three alternative models on the same joint dataset
(Table~\ref{tab:alts}).

\begin{table}[h]
\centering
\small
\setlength{\tabcolsep}{4pt}
\begin{tabular}{lcrrrl}
\toprule
model & $k$ & $\chi^{2}$ & AIC & $\Delta\mathrm{AIC}$ & params \\
\midrule
FREE\_C9                       & 1 & 39.32 & 41.32 & $0.00$  & $\Delta C_{9}=-1.34$ \\
\textbf{VFD\_GREEN\_600CELL}   & 1 & \textbf{37.64} & \textbf{39.64} & $\boldsymbol{-1.67}$ & $A=+1.59$, $\Delta C_{9}^{\mathrm{eff}}=-1.37$ \\
FREE\_C9 + FREE\_C10           & 2 & 39.05 & 43.05 & $+1.73$ & $\Delta C_{9}=-1.40$, $\Delta C_{10}=+0.19$ \\
charm-loop Gaussian            & 2 & 35.41 & 39.41 & $-1.91$ & $A=+1.83$, $\sigma=8.96\,\mathrm{GeV}^{2}$ \\
\bottomrule
\end{tabular}
\caption{Joint-fit comparison on LHCb 2025 four observables.
Adding $C_{10}$ does not help. The free centre-peaked Gaussian
(charm-loop nuisance proxy) achieves a slightly lower AIC than the
kernel at the cost of one extra parameter.}
\label{tab:alts}
\end{table}

\paragraph{Reading.} Adding a free $C_{10}$ shift on top of $C_{9}$
produces a $\chi^{2}$ improvement of only $0.27$, less than the
$+2$ AIC penalty for the extra parameter; the data does not want
$\Delta C_{10}$ at this level. The charm-loop Gaussian — a
two-parameter ansatz with free amplitude and free width centred at
the $J/\psi$\,--\,$\psi(2S)$ midpoint — does fit slightly better
($\Delta\mathrm{AIC}=-1.91$ vs FREE\_C9, vs $-1.67$ for VFD), at the
cost of one extra parameter. Versus VFD it gains $0.24$ in $\chi^{2}$
for the cost of $1$ AIC unit, $\Delta\mathrm{AIC}_{\mathrm{charm\,vs\,VFD}}=-0.24$:
not significant. The fitted width $\sigma\approx 9\,\mathrm{GeV}^{2}$ is
much broader than the natural charm-resonance scale ($\Gamma_{J/\psi}^{2}
\sim 0.1\,\mathrm{GeV}^{2}$), meaning the data is not selecting a
localised charm bump but the same broad centre-peaked shape that the
geometric kernel captures with one parameter rather than two.

\subsection{Form-factor Monte Carlo}

Twenty samples are drawn from the full joint multivariate constraint
on the BSZ $B\to K^{*}$ form-factor coefficients
\citep{Bharucha2016BSZ} via
\texttt{flavio.\allowbreak default\_parameters.\allowbreak get\_random\_all}.
For each draw the BSZ constraints are replaced with delta
distributions at the sampled values;
\texttt{flavio.\allowbreak default\_parameters} is monkey-patched
for the prediction loop and restored on exit. SM and slope tables
are regenerated from scratch per draw (96 \texttt{flavio} calls per
trial).

\begin{table}[h]
\centering
\small
\begin{tabular}{lc}
\toprule
statistic & value \\
\midrule
$n$ samples completed   & $20/20$ \\
mean $A$                & $+1.570$ \\
std $A$                 & $0.163$ \\
$90\,\%$ CI             & $[+1.389,\ +1.801]$ \\
fraction $A<0$          & $0.000$ \\
\bottomrule
\end{tabular}
\caption{BSZ form-factor Monte Carlo statistics. The shared
kernel amplitude shifts by less than $\pm 15\,\%$ under random
form-factor draws and never changes sign.}
\label{tab:ff_mc}
\end{table}

\subsection{Acceptance summary}

\begin{table}[h]
\centering
\small
\begin{tabular}{lp{4.5cm}p{6.0cm}}
\toprule
test (Mode B) & criterion & result \\
\midrule
1. bin bootstrap                & sign-stable, magnitude bounded
                                & \textsc{pass} (Mode B; $99.8\,\%$ positive, CI $\pm 15\,\%$) \\
2. $q^{2}$ region splits        & cross-region consistency, no sign flip
                                & \textsc{pass} (Mode B; 3/3 regions $\Delta\mathrm{AIC}<0$, all $A>0$) \\
3a. FREE\_C9 + FREE\_C10        & adding $C_{10}$ does not displace VFD
                                & \textsc{pass} (Mode B; $\Delta\mathrm{AIC}=+1.73$ vs FREE\_C9) \\
3b. charm-loop Gaussian         & VFD competitive at same shape level
                                & \textsc{partial} (Mode B; $\Delta\mathrm{AIC}_{\mathrm{charm\,vs\,VFD}}=-0.24$, charm $k=2$) \\
4. form-factor MC               & sign-stable, magnitude bounded
                                & \textsc{pass} (Mode B; $100\,\%$ positive, CI $\pm 15\,\%$) \\
5. frozen kernel                & shape never refit
                                & \textsc{pass} (only $A$ fitted everywhere) \\
\midrule
\multicolumn{3}{p{12.0cm}}{\emph{Non-linear stress tests not run;
the linearised $\Delta\mathrm{AIC}$ values quoted above do not
transfer to the non-linear refit.}} \\
\bottomrule
\end{tabular}
\caption{Stress-test acceptance summary on the LHCb 2025
multi-observable joint fit, all under the linearised Mode-B response.
Sign-stability and magnitude-bounded behaviour transfer qualitatively
to the non-linear regime, but Mode-B AIC numbers do not.}
\label{tab:stress_accept}
\end{table}
